\documentclass{article}

\usepackage{graphicx}
\usepackage{minted}

\usepackage[dutch]{babel}

\title{ \includegraphics[scale=1]{gorilla.png} \\ Code Gorilla \\ Use Cases en functioneel ontwerp}
\author{Dagmar Hofman}
\date{}

\begin{document}
	\maketitle
	\newpage
	\tableofcontents
	\newpage
	\section{Inleiding}
	Dit document gaat over \textbf{use cases} en  \textbf{functioneel ontwerp}. Het doel is te \textit{software ontwikkeling in het klein} te leren. Dit heeft de volgende kenmerken: \\
	\begin{itemize}
		\item Het wordt ontwikkeld door \'{e}\'{e}n programmeur
		\item Het kan worden ontwikkeld in een beperkte tijd (uren of dagen maar geen maanden of jaren)
		\item Er is meteen al een duidelijke productomschrijving waaruit nauwkeurig valt op te maken wat het programma doet
		\item Er is geen gebruikersgroep voor wie het programma per se onderhouden moet worden
	\end{itemize}

	\section{Softwareontwikkeling}
	
	\subsection{Voortraject}
	
	In de rest van deze paragraaf gebruiken we als voorbeeld van een groot programma regelmatig een informatiesysteem voor een bedrijf dat verschillende typen containers verhuurt. De huurder neemt de containers op de plaats van vertrek in ontvangst en levert ze weer in op de plaats van bestemming; de verhuurder zorgt ervoor dat ze weer bij een volgende huurder terechtkomen. Het bedrijf heeft twintig kantoren in verschillende landen. Het informatiesysteem moet alle containers gaan volgen en de kantoren in staat stellen om snel te zien of aan een verzoek van een huurder voldaan kan worden en welke containers daar dan het best voor gebruikt kunnen worden. Uiteindelijk moet dat tot een efficiënter gebruik van containers leiden. \\
	Deze cursus gaat niet over softwareontwikkeling in het groot; daar zijn andere cursussen voor. Wat we in deze cursus echter wel willen, is u een programmeerstijl bijbrengen die ook bruikbaar is voor het programmeren in het groot. In een cursus over objectgeoriënteerd programmeren ligt dat ook voor de hand, omdat deze programmeerstijl bij uitstek is ontwikkeld met het oog op grote programma’s. In dat licht moeten de eisen worden gezien die we in deze cursus aan programma’s zullen 	stellen. We zullen daarom in deze paragraaf over de grens van de cursus kijken, naar softwareontwikkeling in het algemeen. \\
	

	Voordat de ontwikkeling start, is er meestal al een heel traject afgelegd, het voortraject. Dit traject start bijvoorbeeld met het signaleren van een probleem of het ontstaan van een nieuwe markt. \\
	Kijk als voorbeeld naar het systeem voor het containerverhuurbedrijf. Het voortraject kan daar begonnen zijn met de constatering dat de	bedrijfsresultaten verslechterden. Toen onderzocht werd hoe dat kwam, bleken de tarieven aan de hoge kant. Een analyse van het bedrijfsproces 	leerde, dat er kosten bespaard konden worden door de containers effici\"{e}nter te gebruiken. Het kwam te vaak voor dat er bijvoorbeeld een stel containers leeg van Rotterdam naar Liverpool werd gebracht, terwijl 	er in Birmingham vergelijkbare containers op een verhuurder stonden te wachten. Die vielen dan weer net onder een ander kantoor, zodat hun beschikbaarheid niet duidelijk werd. Eén informatiesysteem waarop alle kantoren zijn aangesloten, wordt als een mogelijke oplossing gezien. \\
	Als het management voldoende in het idee ziet om er geld in te steken, gaat het project pas echt van start. \\
	Tijdens het voortraject wordt meestal ook het echte ontwikkeltraject opgezet: er wordt bijvoorbeeld bepaald wat het mag kosten en hoe lang het mag duren. \\
	
	\subsection{Analyse}
	
	In de eerste fase, de analyse, moet worden vastgesteld wat er nu eigenlijk precies ontwikkeld moet worden en aan welke eisen het product moet voldoen; pas daarna wordt besloten of het project verder gaat. Een precieze specificatie is nodig om vast te kunnen stellen of het systeem wel met bestaande methoden en technieken gerealiseerd kan worden en of dat kan binnen de randvoorwaarden die in het voortraject zijn bepaald. Ontwikkelaars, opdrachtgevers en toekomstige gebruikers moeten dus rond de tafel gaan zitten om die specificatie op te stellen. Een manier om dat te doen, is het opsporen van alle gewenste gebruiksmogelijkheden (use cases) van het systeem. Iedere gebruiksmogelijkheid is een soort scenario waarin een typische interactie tussen een mogelijke gebruiker en het systeem wordt weergegeven. \\
	
	\includegraphics[scale=1.1]{use\_case\_example.png} \\
	Figuur 1. Voorbeeld van use-cases.
	\newpage
	Twee gebruiksmogelijkheden van het containersysteem zijn bijvoorbeeld de volgende: \\
	
	\begin{itemize}
		\item Het systeem verstrekt op verzoek van een medewerker, wiens taak het
		is om containers toe te wijzen, een overzicht van alle ongebruikte containers (aantal, type en locatie) op een gegeven datum in een gegeven
		gebied (een land, een groep landen of de hele wereld).
		\item Een medewerker voert een mogelijke order in (aantal containers van
		een bepaald type, plaats en datum van vertrek, plaats en datum van
		bestemming) en het systeem wijst op grond daarvan een beschikbare
		verzameling containers aan waarmee zo goedkoop mogelijk (volgens
		een gegeven kostenfunctie) aan de order kan worden voldaan.
		Het is een kwestie van onderhandelen en afwegen welke gebruiksmoge-
		lijkheden in de uiteindelijke specificatie terechtkomen. In het voorbeeld
		kan het realiseren van de tweede gebruiksmogelijkheid zoveel kosten,
		dat het bedrijf besluit de uiteindelijke toewijzing toch met de hand te
		blijven doen en het systeem vooral te gebruiken om alle benodigde in-
		formatie centraal bij te houden en op een overzichtelijke wijze aan alle
		kantoren te leveren.	
	\end{itemize}
	
	Een lijst van gebruiksmogelijkheden vormt de basis van verdere
	gesprekken tussen ontwikkelaar en opdrachtgever.
	
	\subsection{Domeinmodel}
	
	Tijdens analyse (vaak tegelijkertijd met het opstellen van de lijst van
	gebruiksmogelijkheden) wordt ook een domeinmodel opgesteld. In een
	dergelijk model wordt het deel van de werkelijkheid beschreven waar
	het systeem betrekking op heeft. Bij een objectgeoriënteerde analyse
	wordt die beschrijving gegeven in termen van klassen, hun relaties en
	hun interacties. Oppervlakkig lijkt het domeinmodel daarom op een
	ontwerp voor een objectgeoriënteerd programma, maar dat is het niet! \\
	Bij het opstellen van het domeinmodel wordt nog helemaal geen reke-
	ning gehouden met de functionaliteit die de software moet gaan bieden.
	Een domeinmodel is bovendien vaak niet volledig: allerlei details wor-
	den weggelaten om het geheel begrijpelijk en hanteerbaar te houden.
	Ook de essentiële processen uit het domein worden tijdens de analyse
	beschreven, bijvoorbeeld: Hoe wordt een container gevolgd? Hoe wordt
	een order van een klant verwerkt? \\
	
	\subsection{Constructie}
	
	Als het project niet na de analysefase is afgeblazen, kan nu de feitelijke
	constructie van het systeem beginnen. Dat gebeurt vrijwel altijd in fasen.
	De lijst van gebruiksmogelijkheden kan hierbij als uitgangspunt dienen:
	in iedere fase wordt een deel van de gebruiksmogelijkheden gereali-
	seerd. Aan het eind van iedere fase is er dus een werkend systeem, dat
	echter nog niet alles doet wat het volgens de specificaties zou moeten
	doen. Iedere fase in de constructie bestaat zelf weer uit drie stappen:
	ontwerp, implementatie en testen.
	
	\subsection{Ontwerp}
	
	Tijdens het ontwerp van een objectgeoriënteerd systeem wordt vastgesteld uit welke klassen het programma zal gaan bestaan en hoe die
	klassen met elkaar samenhangen. Aan het eind van een ontwerpstap
	moet het volledig duidelijk zijn welke verantwoordelijkheden elke klasse
	in het programma tot dan toe zal hebben. Wat dan nog niet vastligt, is hoe elke klasse haar verantwoordelijkheden zal realiseren. 
	
	\subsection{Ontwerpmodel}
	
	De resulterende klassenstructuur zullen we het ontwerpmodel noemen. Vaak zal het domeinmodel als uitgangspunt dienen voor het opstellen hiervan, maar het ontwerpmodel kan ook aanzienlijk	afwijken van het domeinmodel. Er kunnen extra klassen aan worden
	toegevoegd die niet overeenkomen met elementen uit het domein, maar	die louter een beheertaak hebben. Ook worden er soms andere klassen gekozen, omdat de klassen uit het domeinmodel tot een onhandige of een inefficiënte implementatie zouden leiden. \\
	In deze cursus verdiepen we ons niet in zulke gevallen, maar het is goed om u te realiseren dat bij programmeren in het groot het modelleren van het domein en het ontwerpen van de software gescheiden activiteiten zijn die tot verschillende resultaten kunnen leiden. \\
	
	\subsection{Technisch ontwerp}
	
	Voor het softwareontwerp wordt daarom ook wel de term technisch ontwerp gebruikt.
	Verder moet u zich realiseren dat een ontwerp vaak een uitbreiding is
	van een vorig ontwerp. De constructie verloopt immers in fasen; in elke
	fase worden gebruiksmogelijkheden toegevoegd. Soms kan daarvoor
	worden volstaan met het uitbreiden van bestaande klassen, soms moeten
	nieuwe klassen worden toegevoegd.
	
	\subsection{Implementatie}
	
	Als er een bevredigend ontwerp ligt, kan begonnen worden met de
	implementatie. Per klasse moet worden vastgesteld hoe deze klasse haar
	verantwoordelijkheden gaat verwezenlijken. Eventueel kan dit worden
	vastgelegd in een apart implementatiemodel waarin alle attributen met
	hun typen zijn vastgelegd en alle methoden met hun signatuur (de kop
	van de methode) en met een beschrijving van hun werking. \\
	Vervolgens wordt de klasse gecodeerd. Omdat, tijdens de ontwerpfase,
	de interface van een klasse gedetailleerd is vastgelegd, kunnen in deze
	fase verschillende programmeurs onafhankelijk van elkaar aan verschil-
	lende klassen werken.
	
	\subsection{Testen}
	
	Het testen van een bepaalde constructiestap verloopt ook weer in stappen.
	Elke klasse wordt eerst afzonderlijk getest. De programmeur zal daarvoor
	een kleine testomgeving construeren. In Objectgeoriënteerd programmeren maakten we hiervoor gebruik van JUnit. Als alle klassen lijken te
	werken, worden ze samengevoegd en wordt het programma als geheel
	getest. Hierbij zullen eventuele misverstanden tussen de verschillende
	programmeurs aan het licht komen. Als het implementatiemodel volledig
	en eenduidig is en alle programmeurs volmaakt zijn, dan zijn die misver-
	standen er niet. De praktijk leert dat aan die voorwaarden vaak niet is
	voldaan. \\
	Iedereen die regelmatig met computers werkt, weet hoe moeilijk testen
	is: in ieder softwarepakket van enige omvang blijken toch altijd weer
	fouten te zitten. Om dat aantal zo klein mogelijk te houden, is een goede
	teststrategie van zeer groot belang. Het zal duidelijk zijn dat voor ieder
	programma, hoe klein ook, slechts een miniem deel van alle mogelijke
	combinaties van invoerwaarden daadwerkelijk kan worden uitgeprobeerd. De kunst is nu dat deel zo te kiezen, dat het zo veel mogelijk
	representatief is voor alle invoer.
	
	\subsection{Afronding}
	
	Er komt een moment waarop ook de laatste gebruiksmogelijkheid is
	toegevoegd en getest en het systeem dus in principe af is. Er is dan vaak
	nog een afrondingsfase waarin de opdrachtgever bijvoorbeeld het systeem
	al ter beschikking krijgt om het uit te proberen. Of, als het een nieuw pak-
	ket betreft, wordt op dat moment misschien een gratis te downloaden
	bètaversie op internet gezet, zodat potentiële klanten deze uit kunnen
	proberen. Daar komen altijd nog fouten uit, die in de afrondingsfase
	verbeterd kunnen worden. Tot slot wordt het systeem dan echt overgedragen of vrijgegeven voor verkoop. \\
	Het hier geschetste ontwikkeltraject is niet het enig mogelijke. Er zijn
	bijvoorbeeld andere faseringen mogelijk. Als ook de analyse en de
	afronding in fasen worden uitgevoerd, bestaat het hele traject uit het
	herhaaldelijk (iteratief) doorlopen van de stappen analyse, ontwerp,
	implementatie, testen en afronding van de huidige fase. De opdracht-
	gever heeft dan al snel een klein werkend systeem ter beschikking.
	Na iedere volgende fase zullen de mogelijkheden zijn uitgebreid.
	Deze iteratieve aanpak is te beschouwen als een eenvoudige vorm van
	het Unified Process (UP). Kenmerkend voor het UP is het iteratief en
	incrementeel ontwerpen en implementeren van een systeem.
	
	\subsection{Prototype}
	
	Op deze wijze kan ook een prototype worden ontwikkeld: een experimen-
	tele versie van een te realiseren systeem. Meestal wordt een prototype
	gebruikt om inzicht te krijgen in de werking van een uiteindelijk programma. Aan de hand van zo’n prototype kan worden getoetst of aan
	alle eisen is voldaan en of het programma aan de verwachtingen van de
	opdrachtgever voldoet. Doordat de opdrachtgever in een vroeg stadium
	met het prototype kan ‘spelen’, is het mogelijk het ontwerp tijdig in een
	volgende fase bij te stellen. \\
	
	Aan de andere kant zou het volledige ontwerp en de implementatie in
	een keer kunnen worden gedaan, zodat de constructie uit slechts één
	stap bestaat. Hoe groter het systeem is, hoe sterker dit laatste echter moet
	worden afgeraden: behoorlijk testen van heel veel code ineens is vrijwel
	onmogelijk. Op de voor- en nadelen van verschillende faseringen gaan
	we hier verder niet in. \\
	
	\subsection{Onderhoud}
	
	Als het systeem is opgeleverd of op de markt is gebracht, is daarmee
	de kous nog lang niet af. Het systeem moet namelijk ook onderhouden
	worden en in de praktijk is daarmee vaak veel meer tijd en geld gemoeid
	dan met de oorspronkelijke constructie. \\
	Waar bestaat dat onderhoud uit? Ten eerste zullen er, hoe goed er ook
	getest is, nog steeds fouten in het programma zitten, die in een volgende
	versie verbeterd moeten worden. Ten tweede zullen er veranderingen in
	het domein optreden of in de omgeving waarin het programma draait,
	die aanpassingen in het systeem nodig maken. Ten derde kunnen er uit-
	breidingen van het programma nodig zijn, wat wil zeggen dat er nieuwe
	gebruiksmogelijkheden aan moeten worden toegevoegd. \\
	Ook hiervan kan de oorzaak liggen in veranderingen in het domein,
	maar het kan ook zijn dat het regelmatige gebruik van het systeem de
	gebruiker op nieuwe ideeën brengt: ‘het zou toch ook wel handig zijn
	als ...’, of dat de voortdurende voortschrijdende techniek nieuwe
	mogelijkheden binnen bereik brengt. \\
	\newpage
	Een programma kan op deze wijze vele jaren meegaan, in steeds weer
	nieuwe gedaanten, zonder ooit echt uit de roulatie te worden genomen.
	Uit het oogpunt van softwarearchitectuur zou het misschien verstandig
	zijn om programma’s elke vijf of hoogstens tien jaar te vervangen door
	een volledig nieuwe versie die van de grond af aan is herontwikkeld.
	Daar zijn echter vaak zulke hoge kosten mee gemoeid dat een dergelijke
	oplossing economisch niet haalbaar is. Het millenniumprobleem, waar
	in de laatste jaren van de vorige eeuw veel aandacht voor was, kan hier
	gebruikt worden als illustratie: de oerversie van sommige programma’s
	waarin jaartallen met slechts twee cijfers gerepresenteerd werden,
	dateerde nog uit de jaren zestig van de vorige eeuw! \\
	In de totale levensloop van een groot systeem is onderhoud dus
	minstens zo belangrijke factor als de oorspronkelijke ontwikkeling. \\
	\newpage
	\section{Wat is een goed programma}
	
	Welke eisen kunnen we nu, in het licht van de zojuist geschetste levensloop van softwaresystemen, stellen aan ontwerp en implementatie van
	een dergelijk systeem?
	
	\subsection{Correctheid}
	
	Allereerst moet een programma uiteraard correct zijn: het moet doen
	wat het volgens de specificaties zou moeten doen. Als gebruik gemaakt
	wordt van formele specificatietechnieken, kan de correctheid van een
	programma in principe met wiskundige methoden bewezen worden. \\
	In deze cursus zullen we dat niet doen; in plaats daarvan zullen we, net
	als vrijwel altijd in de praktijk gebeurt, via testen fouten opsporen en die
	vervolgens herstellen. Realiseert u zich wel dat deze strategie beperkingen
	kent. Een goede teststrategie kan veel fouten aan het licht brengen, maar
	via testen kan nooit worden aangetoond dat het uiteindelijke programma
	honderd procent correct is. \\
	Correctheid is overigens een vrij beperkt begrip. In feite moet een
	programma doen wat de opdrachtgever wil of wat de toekomstige
	gebruikers verlangen. Als de specificatie daarmee achteraf niet in
	overeenstemming blijkt en de gebruikers ervaren het programma als
	onhandig of zelfs onbruikbaar, dan is het programma niet goed, zelfs
	niet wanneer bewezen is dat het correct is. \\
	\newpage
	\subsection{Robuustheid}
	
	Correctheid alleen is bovendien niet genoeg. Een programma moet ook
	robuust zijn, ofwel: het moet tegen mogelijke fouten bestand zijn. Dat
	kunnen fouten van de gebruiker zijn: bijvoorbeeld het invoeren van let-
	ters als een getal wordt verwacht, of het opgeven van een niet-bestaande
	datum, of het opvragen van gegevens uit een bestand die daar niet in
	zitten, of het kiezen van een verkeerde menuoptie waardoor een onbedoelde interactie wordt gestart of juist het abusievelijk afbreken van een
	lopende interactie. \\
	In zulke gevallen verwachten we een redelijke respons van de software
	waarmee we werken: een geluid dat aangeeft dat we iets doen wat niet
	klopt, een heldere maar beknopte foutmelding, een mogelijkheid om de
	gestarte interactie meteen weer af te breken, of een waarschuwing dat
	we op het punt staan (veel) werk weg te gooien. Het programma moet
	bovendien bestand zijn tegen allerlei andere onvoorziene omstandig-
	heden, zoals een harde schijf die vol is zodat er geen data kunnen
	worden weggeschreven. \\
	
	\subsection{Begrijpelijke code}
	
	Wil een programma onderhouden kunnen worden, dan moet de code
	begrijpelijk zijn. De programmeur die een fout moet verbeteren of een
	wijziging aanbrengen, is vaak een andere dan degene die de code oor-
	spronkelijk heeft geschreven. Als de tijd tussen schrijven en veranderen
	langer dan een paar maanden is, maakt dat bovendien weinig verschil
	meer: programmeurs staan dan net zo vreemd tegenover hun eigen code
	als tegenover die van ieder ander. \\
	Een manier om de begrijpelijkheid van code te verbeteren, is een goede
	documentatie van het programma in de vorm van uitgebreid commentaar. In onze Java-cursussen gebruiken we daarvoor veelal javadoc,
	aangevuld met gewoon commentaar bij lastige stukken broncode. \\
	
	\subsection{Uitbreidbaarheid}
	
	Ook de structuur van het programma moet dusdanig zijn, dat het makke-
	lijk te wijzigen en uit te breiden is. Een wijziging is makkelijker naarmate
	deze op minder delen van het programma invloed heeft. Wanneer op
	twintig plekken iets gewijzigd moet worden, zien we er gauw een over
	het hoofd. Een uitbreiding is gemakkelijker naarmate het nieuwe stuk
	zelfstandiger te ontwikkelen is en tot minder wijzigingen elders in het
	programma leidt. \\
	Er zijn nog meer eisen waar een programma aan moet voldoen, maar
	daar besteden we in deze cursus minder aandacht aan. We noemen er
	nog twee. \\
	\subsection{Herbruikbaarheid}
	
	Een gewenste eigenschap die objectori\"{e}ntatie populair heeft gemaakt,
	is herbruikbaarheid van delen van het systeem. Soms kan het veel ont-
	wikkeltijd besparen als een klasse die voor een bepaalde toepassing
	gemaakt is, ook in een andere toepassing kan worden ingezet. Het
	belang van deze eigenschap is u in feite al bekend: de Java API is in
	wezen niets anders dan een verzameling herbruikbare klassen	
	
	\subsection{Effici\"{e}ntie}
	
	Tot slot is voor sommige programma’s efficiëntie heel belangrijk. Dit geldt
	bijvoorbeeld voor:
	\begin{itemize}
		\item toepassingen waarbij het programmaverloop een proces in de buiten-
		wereld moet bijhouden, zoals een industriële robot die moet reageren op
		aanvoer van onderdelen op een band
		\item toepassingen waarbij een factor tien in geheugengebruik of tijd net de
		grens tussen bruikbaar en onbruikbaar vormt: een weervoorspelling voor
		de komende week mag een dag rekentijd vragen, maar geen tien dagen
		\item pakketten waarbij de interactie met de gebruiker zeer intensief is en
		wachten al snel hinderlijk wordt: besturingssystemen, tekstverwerkers,
		ontwikkelomgevingen.
	\end{itemize}

	Al deze eisen zijn op zich redelijk voor de hand liggend. De vraag is
	echter, hoe we die eisen vertalen in eigenschappen van programmacode.
	Daarbij zullen we ons in de rest van deze paragraaf concentreren op
	begrijpelijkheid en gemak van wijzigen en uitbreiden. In tegenstelling
	tot bijvoorbeeld robuustheid en efficiëntie zijn dit namelijk de eisen
	waarmee van het begin af aan rekening moet worden gehouden. \\
	Het is mogelijk om een niet-robuust, maar begrijpelijk en makkelijk
	te wijzigen programma, achteraf alsnog robuust te maken. Een on-
	doorgrondelijk en moeilijk te wijzigen programma achteraf alsnog
	begrijpelijk maken is niet echt onmogelijk, maar wel uiterst moeilijk.
	Vaak vereist het een volledig nieuw ontwerp en implementatie. \\
	
	\subsection{Kleine klassen en methoden}
	Een heel belangrijke eigenschap die de begrijpelijkheid van code bevordert is eenvoud. In het algemeen zullen klassen en methoden moeilijker te
	begrijpen zijn naarmate ze groter zijn. Het is daarom goed om te streven
	naar kleine klassen en methoden. Zoals u bij het gebruik van de API specification gemerkt zult hebben, voldoen de klassen uit de Java API niet altijd
	aan deze norm. 
	\subsection{Eenvoudige control flow}
	Ook is een methode begrijpelijker naarmate de control flow eenvoudiger is. Daarom proberen we een diepe nesting van while’s, for’s en if’s te vermijden. en klein maar vaak gezien voorbeeld van moeilijk te lezen code is het volgende: \\
	
	\begin{minted}{java}
	if (!test) {
		return false;
	} else {
		return true;
	}
	\end{minted}
	Dit kan korter en daarmee duidelijker opgeschreven worden:

	\begin{minted}{java}
	return test;
	\end{minted}
	
	Een andere belangrijke manier om eenvoud te bereiken, is verwoord
	in het volgende principe: ieder onderdeel van de code dient slechts één
	verantwoordelijkheid te hebben. Dit principe geldt voor alle niveaus:
	opdrachten, methoden en klassen. \\
	
	\subsection{Lokaliteit}
	
	Op methodeniveau proberen we methoden te vermijden die zowel een waarde opleveren als een of meer attributen wijzigen. Een dergelijke methode heeft twee verantwoordelijkheden: iets aan de toestand van het	object veranderen en een resultaat teruggeven.
	Er zijn op deze regel overigens wel uitzonderingen. Het is bijvoorbeeld	niet verkeerd om een methode een waarde terug te laten geven die laat zien of de actie die de methode moet uitvoeren, geslaagd is. Een voorbeeld hiervan is de methode public boolean add(E e) in de klasse ArrayList. Deze methode geeft de waarde true terug als het toevoegen van het element geleid heeft tot een verandering van de arraylist, en false als het toevoegen niet gelukt is. Op klassenniveau zullen we het principe van gescheiden verantwoordelijkheden als een van de leidraden gebruiken bij de keuze van klassen in een ontwerp. \\ 
	Het is een goede gewoonte om bij het ontwerp het doel van een klasse – we kunnen ook zeggen: de verantwoordelijkheid van de klasse – in een kort zinnetje expliciet te vermelden. Als in zo’n zinnetje het woordje ‘en’ voorkomt, moet op zijn minst onderzocht worden of die klasse niet beter gesplitst kan worden, zoals in: deze klasse representeert een order en co\"{o}rdineert de toewijzing van containers aan die order. Het antwoord is in elk geval ‘ja’ als dat kan gebeuren zonder dat er veel extra interactie tussen die klassen nodig is. Gescheiden verantwoordelijkheden maken code gemakkelijker te begrijpen en alleen al daarom ook gemakkelijker te wijzigen en uit te breiden. Een wijziging of uitbreiding is echter ook des te gemakkelijker, naarmate deze op minder plaatsen in de code van invloed is. Stel bijvoorbeeld dat een wijziging invloed heeft op de signatuur van drie methoden uit drie verschillende klassen, en dat elk van die drie methoden op vijf verschillende plaatsen wordt aangeroepen. Om die ene wijziging door te voeren, moeten we dan op tenminste achttien plaatsen in de code iets veranderen: de drie methoden plus de vijftien aanroepen. Niet alleen moet elk onderdeel van de code slechts \'{e}\'{e}n verantwoordelijkheid hebben, iedere verantwoordelijkheid van het systeem als geheel moet bij voorkeur ook slechts op één plek in de code gerealiseerd zijn. Als we dan iets wijzigen aan die verantwoordelijkheid, dan hoeven we alleen op die plek iets te veranderen. We noemen deze eis aan de code het principe van lokaliteit.
	
	\subsection{Constanten}
	
	Ook dit principe speelt op verschillende niveaus. De eis van lokaliteit
	is bijvoorbeeld een van de redenen dat we constanten definiëren. Stel er
	zijn in het containersysteem vijf verschillende typen containers en dat
	aantal komt op tien plaatsen in de code voor. Als we op al die tien
	plaatsen het getal 5 neerzetten en er komt een type container bij, dan
	moeten we tien wijzigingen aanbrengen. Is er ergens in het systeem een
	constante AANTALCONTAINERTYPEN gedefinieerd, dan hoeven we alleen
	de waarde van die constante te veranderen. \\
	
	Ook onze voorkeur voor attributen die van buiten de klasse niet ge-
	wijzigd kunnen worden (information hiding), is terug te voeren op het
	principe van lokaliteit: de verantwoordelijkheid voor het beheer van een
	dergelijk attribuut ligt geheel binnen de klasse, in plaats van verspreid
	door het hele programma. Met het oog daarop, maken we attributen
	vrijwel altijd private en zijn we voorzichtig met het exporteren van
	referenties naar attributen van een objecttype.
	Het bereiken van lokaliteit is in de praktijk een van de moeilijkste zaken
	bij objectgeoriënteerd programmeren (en bij programmeren in het algemeen). \\
	
	Op grond van het principe van lokaliteit willen we aan een klasse
	gemakkelijk een nieuwe subklasse kunnen toevoegen. Stel bijvoorbeeld
	dat er in het containersysteem een klasse Container is, met een subklasse
	voor elke soort container. Als we nu een nieuw soort container willen
	toevoegen, zullen we in elk geval een nieuwe subklasse van Container
	moeten definiëren. Daarnaast zijn wijzigingen nodig op plekken in de
	code waar instanties van de subklassen van Container gecreëerd worden
	en op plekken waar het nodig is na te gaan met welk type container we
	nu precies te maken hebben. Lokaliteit vereist dat dit op zo min mogelijk
	plaatsen voorkomt. In de volgende leereenheid zullen we zien hoe Java
	ons daarbij helpt. \\
	
	\subsection{Lege koppeling}
		
	Nog een manier die we hier willen noemen om een programma
	eenvoudiger te maken, is het aantal associaties tussen klassen te
	beperken ofwel de koppeling tussen de klassen laag te houden. Als we een
	systeem hebben met tien klassen en iedere klasse is afhankelijk van alle
	andere, dan is dat systeem vrijwel zeker moeilijk te doorgronden en te
	wijzigen. Een wijziging in een klasse maakt dan al snel wijzigingen in
	alle andere klassen noodzakelijk. Hergebruik van één klasse in een ander
	systeem is al helemaal onmogelijk.
	
	\subsection{Objectkeuze}
	
	Het is lang niet gemakkelijk om programmacode te schrijven met al de
	gewenste eigenschappen. Het kiezen van goede namen is vooral een
	kwestie van zelfdiscipline, evenals het schrijven van commentaar. Maar
	voor het schrijven van code die ook aan de andere eisen voldoet, is
	ervaring nodig. Een beginnend programmeur ziet vaak maar één manier
	om een probleem aan te pakken. Hij zal daarom al gauw tot de conclusie
	komen dat déze grote klasse echt niet gesplitst kan worden, of dat déze
	methode echt zowel een waarde moet opleveren als de toestand van een
	object moet veranderen. Soms heeft een andere oplossing inderdaad
	meer nadelen dan voordelen, maar meestal ziet een programmeur met
	enige ervaring wel hoe het anders kan. Nog veel meer ervaring is nodig
	om klassen dusdanig te kiezen en uit te werken, dat de resulterende code
	voldoet aan onder andere de principes van gescheiden
	verantwoordelijkheden, lage koppeling en lokaliteit. \\
	De verschillende eisen aan de code zijn bovendien soms met elkaar in
	conflict. Soms gaan een vergaande scheiding van verantwoordelijkheden
	en een grote lokaliteit ten koste van de eenvoud van een ontwerp.
	
\end{document}          
